\section{Synthèse Exécutive / Executive Summary}

% ─── VERSION FRANÇAISE ───────────────────────────────────────────────
\subsection*{Synthèse exécutive (Français)}

\begin{guideline}
Cette synthèse doit pouvoir « vivre seule », sans le reste du rapport.
Elle est destinée à un lecteur non-initié : soyez percutant, synthétique
et pédagogique. Elle couvre : le contexte de l'entreprise, la
problématique adressée, votre démarche, vos réalisations clés et un
bilan. Environ 1 page.
\end{guideline}

\instruction{[Rédigez ici votre synthèse en français. Structurez-la
autour de 4 à 5 paragraphes courts : contexte, problématique, démarche,
réalisations, bilan. Évitez le jargon technique brut — un lecteur non
spécialiste doit comprendre l'essentiel.]}

\vspace{1cm}

% ─── VERSION ANGLAISE ────────────────────────────────────────────────
\subsection*{Executive Summary (English)}

\begin{guideline}
This summary must stand alone, without the rest of the report. It is
aimed at a non-specialist reader: be impactful, concise and clear. Cover
the company context, the problem addressed, your approach, key
achievements and an overall assessment. Approximately 1 page.
\end{guideline}

\instruction{[Write your executive summary here in English. Structure it
around 4 to 5 short paragraphs: context, problem statement, approach,
achievements, conclusion. Avoid raw technical jargon — a non-specialist
reader should grasp the essentials.]}